\section{Related work}

\paragraph{Limits of reference-based metrics} have been explored for text summarization tasks and different solutions have been proposed to address them.
\citet{scialom-etal-2019-answers, scialom-etal-2021-questeval, fabbri-etal-2022-qafacteval} use question generation pipelines to evaluate the faithfulness of summaries without relying on references.
\citet{spice} study the use of intermediate dependency parse trees to rate natural image captions. 
For semi-structured data such as table-to-summary generation, PARENT~\cite{dhingra-etal-2019-handling} demonstrates the limitation of BLEU as it do not highlight the key knowledge from the table, whereas \citet{opitz-frank-2021-towards} focuses on generating text from abstract meaning representations. 
\citet{gehrmann2022tata} observed the poor correlation of BLEURT-20 to human ratings and proposed STATA, a learned metric using human annotation.

In this work we explore building a reference-free metric for chart summary that does not require human-annotated references. We show that our metric \metric has much higher correlation with human judgment than reference-based metrics such as BLEU.

\paragraph{Chart-to-summary generation} has recently gained relevance within multimodal NLP.
\citet{obeid-hoque-2020-chart} created the Chart-to-Text dataset, using charts extracted from Statista. \citet{kantharaj-etal-2022-chart} extended it with more examples from Statista and Pew Research.
Besides efforts on evaluation, multiple modeling methods have been proposed to reduce hallucinations. The approaches can be roughly divided into (1) pipeline-based methods which first extract chart components (e.g. data, title, axis, etc.) using OCR then leverage text-based models to further summarize the extracted information \citet{kantharaj-etal-2022-chart, choi2019visualizing}; (2) end-to-end models which directly input chart-attribute embeddings to Transformer-based models for enabling structured understanding of charts \cite{obeid-hoque-2020-chart}. 

In this work we explored both (1) and (2). The best approach \pipeline{} generally follows the idea of (1). Instead of relying on OCR we use a de-rendering model for extracting structured information in charts and we explore a self-critiquing pipeline with LLMs for the best quality chart summarization. 
